\documentclass[11pt]{article}
\usepackage[margin=1in]{geometry}
\usepackage{fontspec}
\setmainfont{Arial}
\usepackage{graphicx}
\usepackage{fvextra}
\usepackage{enumitem}
\setlist{nosep}
\setlength{\parindent}{0pt}
\setlength{\parskip}{0.55em}
\begin{document}

\begin{center}\LARGE\textbf{The second unified View regression: two probes, one table, one local consumer}\end{center}
\section*{1 · QA inputs}
The incoming side: this View binds two of QBt1's three answered Probes, reused by path into its own input/QA-probes/ copy.

\begin{Verbatim}[fontsize=\small]
QA-bank ─▶ answered Probe (definition) ─┐
QA-bank ─▶ answered Probe (boundary)   ─┴─▶ QY3 input collection
\end{Verbatim}
This View imports two answered QA Probes: the canonical trait definition and the measurement boundary.

\begin{quote}\footnotesize\textbf{Evidence card.} Card two answered QA Probes: QI1 · QA input · Bindings: views/QY3-for-view/input/QA-probes/1-canonical-definition.md, views/QY3-for-view/input/QA-probes/3-measurement-boundary.md. Status: answered. Role: supply the material organized by this View.\end{quote}
The observable-signal Probe is deliberately excluded from this regression; only construct meaning and measurement boundary carry through to the Display.

\section*{2 · View body}
The readable middle: two short subsections hold the actual substance, while Cards attach provenance to exact words.

\begin{Verbatim}[fontsize=\small]
readable sentence + exact-span Card ─▶ inspect evidence without leaving the body
\end{Verbatim}
\subsection*{2.1 · Trait meaning}
Agreeableness is an interpersonal trait expressed through warmth and cooperation.

\begin{quote}\footnotesize\textbf{Evidence card.} Card interpersonal trait: EC1 · Literature · Finding: agreeableness is treated as an interpersonal construct. Bearing: supports. Boundary: trait description only. Binding: views/QY3-for-view/input/QA-probes/1-canonical-definition.md. Freshness: current. Used by: QY3-Display1.\end{quote}
The literature-side construct frame resolves through the real bibliography (John et al., 1999).

\subsection*{2.2 · Measurement boundary}
The score is not treated as an error-free measure of latent personality.

\begin{quote}\footnotesize\textbf{Evidence card.} Card not treated as an error-free measure of latent personality: EC2 · Construct · Finding: the current evidence does not justify direct latent-personality measurement. Bearing: bounds this specimen. Boundary: patient-perceived manifestation only. Binding: views/QY3-for-view/input/QA-probes/3-measurement-boundary.md. Freshness: current. Used by: QY3-Display1.\end{quote}
The body stays a human-readable synthesis; the two Evidence Cards let a reader inspect why each exact phrase is present.

\section*{3 · Displays}
The visible output: QY3-Display1 is the trait-boundary table. Its current raster is embedded below; the same preview appears first in its Display Card, before its bindings, files, and acceptance state.

\begin{center}\includegraphics[width=0.94\linewidth]{assets/display-1.png}\end{center}
\begin{quote}\footnotesize\textbf{Evidence card.} Card QY3-Display1: QY3-Display1 · Display · Binding: views/QY3-for-view/output/QY3-Display1-trait-boundary-table/output.md. Kind: table. Reader job: separate construct meaning from the measurement boundary the two bound Probes support. View-body bindings: EC1, EC2. Status: rendered. Acceptance: waiting. Files: views/QY3-for-view/output/QY3-Display1-trait-boundary-table/preview.png, views/QY3-for-view/output/QY3-Display1-trait-boundary-table/preview.pdf, views/QY3-for-view/output/QY3-Display1-trait-boundary-table/float.tex.\end{quote}
preview.pdf is the printable inspection surface for the same table; float.tex is the printable wrapper the review build embeds.

\section*{4 · Consumers}
The outgoing side: this regression hands off to one local consumer Page instead of a shared QS-consumers Page, so the run stays self-contained.

\begin{Verbatim}[fontsize=\small]
View ─▶ QY3-Display1 consumer + placement + gate state
\end{Verbatim}
The local boundary note plans to consume QY3-Display1 for its construct-boundary passage.

\begin{quote}\footnotesize\textbf{Evidence card.} Card local boundary note: C1 · Consumer · Target Page: QY3-Consumer-boundary-note · Binding: consumers/QY3-Consumer-boundary-note.md. Uses: QY3-Display1. Placement: Local regression consumer / construct-boundary passage. Status: planned, blocked on View and Display acceptance.\end{quote}
\section{References}\begin{itemize}
\item John, Oliver P.; Srivastava, Sanjay (1999). The Big-Five Trait Taxonomy: History, Measurement, and Theoretical Perspectives. Handbook of Personality: Theory and Research
\end{itemize}
\end{document}
